% Conclusion (issue #5; final report #63). Every number is measured, from results/phase2/
% (#69, #64) and results/phase2/synth/synth_summary.csv (#75).

\section{Conclusion}
This project built a parameterized output-stationary systolic array for fixed-point matrix multiplication in synthesizable Verilog-2005, and took it from a single $N \times N$ product in Phase~1 to products of any size in Phase~2. Two small changes outside the array made this possible: a run-time inner length, which sets how long the skew buffers feed operands and how many steps the controller runs, and an accumulate start that skips the clear, so the partial sums of a tile stay in their PEs from one chunk of the inner dimension to the next. The array and its PEs did not change. Every measured run meets the written timing contract of $k_{\mathrm{len}} + 2N$ busy cycles, and all 26 tiled products, up to $9 \times 33 \times 7$ and $4 \times 64 \times 4$, match the golden model. A sequential baseline with the same interface lets the same testbenches measure the benefit of the array: 5.3 times fewer busy cycles for $k_{\mathrm{len}} = 4$, rising to 10.7 for $k_{\mathrm{len}} = 16$ and 11.0 on the longest tiled product, because the fill and drain of the wavefront are paid once per run. On an ECP5 LFE5U-25F the $4 \times 4$ array uses 16 multiplier blocks, 2506 LUTs and 2287 flip-flops and reaches 56.83~MHz, against 52.16~MHz for the baseline, so a run with $k_{\mathrm{len}} = 16$ finishes 11.6 times sooner. The parallelism costs no clock frequency and no extra LUTs: the array uses 33\,\% fewer LUTs than the baseline and 1.4 times its flip-flops, and its speed comes from its 16 multipliers.

The results also show the limits of the design. The host writes and reads one element per cycle, so loading a run's operands takes $N k_{\mathrm{len}}$ cycles, more than the run itself, and the array is busy for at most 30\,\% of a tiled product: the end-to-end speed-up of an $8 \times 8 \times 8$ product is 2.7 against 8.0 for the array alone. The wavefront's $2N - 2$ cycles of fill and drain are paid on every run, so a PE is busy for at most 72.7\,\% of the compute steps, and the edge tiles of a product whose size is not a multiple of $N$ spend as long on zero padding as on useful work. No configuration meets the 100~MHz target, because the combinational skew lanes feed the multiply-accumulate directly, and an $8 \times 8$ array needs more multipliers than the device has. Results stay at 32 bits, with no rounding or saturation back to 8 bits. The design was evaluated in simulation and by static timing analysis only, since no FPGA board was available. These limits set the next steps: wider or double-buffered operand ports that overlap loading with computing, registered skew lanes for a faster clock, rounding and saturation of the results, and running the design on an FPGA board.
